(1) 设切向静摩擦 \(f\) 沿下坡为正，则质心方程 \(ma=mg\sin\theta+F+f\)。接触摩擦对质心向下滚动角方向产生反向力矩，因此 \(I\alpha=\tau-fR\)。无滑动运动学约束是 \(a=R\alpha\)，这三个方程必须联立；若先假定摩擦总向上，可能错误处理较强驱动力矩情形。

(2) 由转动方程得 \(f=\tau/R-Ia/R^2\)。代入平动方程，\((m+I/R^2)a=mg\sin\theta+F+\tau/R\)，因此 \(a=(mg\sin\theta+F+\tau/R)/(m+I/R^2)\)。再代回得到 \(f=\tau/R-(I/R^2)(mg\sin\theta+F+\tau/R)/(m+I/R^2)\)。\(f\) 的正负由参数决定，且 \(\tau=mR^2a/I\) 时可为零。

(3) 法向反力 \(N=mg\cos\theta\)，持续纯滚动必要条件 \(|f|\le\mu_s mg\cos\theta\)；对于单一库仑静摩擦上界和给定恒定参数，该条件亦充分。把上述 \(f\) 表达式代入得完全显式的不等式。若它被违反，则必须重写滑动摩擦模型，不能在旧的 \(a=R\alpha\) 公式上机械取 \(\mu_kN\)。

(4) 从静止出发，\(v^2=2as,\ \omega=v/R\)，因此总动能 \(\Delta K=(m+I/R^2)v^2/2=(mg\sin\theta+F+\tau/R)s\)。右端正是重力功 \(mg s\sin\theta\)、推力功 \(Fs\) 与外力矩功 \(\tau(s/R)\) 之和。接触点瞬时速度为零，静摩擦的质心平动功与对质心的转动功相互抵消，故理想纯滚动不发生摩擦热耗散。
